\documentclass[11pt,twoside]{article}   % For Latex2e
\usepackage{amsthm, amsmath}   % For Latex2e

\usepackage[]{hyperref}

% SI NO SE CARGA hyperref, INCLUIR ESTO:
%% \usepackage{url}

% SI NO SE CARGA NI url NI hyperref, INCLUIR ESTO:
%\def\url#1{{\def~{\textasciitilde}\tt#1}}

% SI NO SE CARGA hyperref, INCLUIR ESTO:
%\def\href#1#2{\url{#2}}

\pagestyle{myheadings}
\markboth{\footnotesize Actas EACA--2001.
Universidad de La Rioja}{\footnotesize M. Benito, J. L. Varona,
``Aliquot sequences starting with a number under 10000''}

\setcounter{page}{90} %%%%% NUMERO DEFINITIVO EN LAS ACTAS DEL CONGRESO
 \textwidth=6in \hoffset=-1.3cm \textheight=20.5cm
 \oddsidemargin 1.5cm \evensidemargin 1.5cm

\theoremstyle{plain}
% en castellano :
\newtheorem{teorema}{Teorema}[section]
\newtheorem{lema}[teorema]{Lema}
\newtheorem{proposicion}[teorema]{Proposici\'on}
\newtheorem{corolario}[teorema]{Corolario}
% in english :
\newtheorem{theorem}{Theorem}[section]
\newtheorem{lemma}[theorem]{Lemma}
\newtheorem{proposition}[theorem]{Proposition}
\newtheorem{corollary}[theorem]{Corollary}

\theoremstyle{definition}
% en castellano :
\newtheorem{definicion}[teorema]{Definici\'on}
\newtheorem{ejemplo}[teorema]{Ejemplo}
\newtheorem{algoritmo}[teorema]{Algoritmo}
% in english :
\newtheorem{definition}[theorem]{Definition}
\newtheorem{example}[theorem]{Example}
\newtheorem{algorithm}[theorem]{Algorithm}

\theoremstyle{remark}
% en castellano :
\newtheorem{nota}[teorema]{Nota}
% in english :
\newtheorem{remark}[theorem]{Remark}

\newenvironment{dem}{\noindent{\it Demostraci\'on}.}{\qed}

%%%%%%%%%%%%%%%%%%%%%%%%%%%%%%%%%%%%%%%%%%%%%%%%%%%%%%%%%%%%%
%%%%%%%% TO BUILD THE TABLE:
%%%%%%%%%%%%%%%%%%%%%%%%%%%%%%%%%%%%%%%%%%%%%%%%%%%%%%%%%%%%%
\newcommand{\info}{\(\begin{array}{@{}c@{}}{n\strut}\\{k}\\%
    {\text{digits}}\\{\text{guide}}\end{array}\)}
\newcommand{\num}[4]{\(\begin{array}{@{}c@{}}{#1\strut}\\{#2}\\%
    {#3}\\{#4\strut}\end{array}\)}
\newcommand{\pp}{\!\cdot\!}
\newcommand{\nodriver}{{}^{(*)}}                 % WITHOUT COLOR
%\newcommand{\nodriver}{{}^{\Magenta(*)\Black}}  % WITH COLOR in Textures 1.7
%\newcommand{\nodriver}{{}^{\textcolor{magenta}{(*)}}}  % WITH COLOR and package "color"
%\newcommand{\nodriver}{}                        % NOTHING
%\newcommand{\ldw}[1]{\mbox{\scriptsize\ (#1)}}  % Last driver with [1] digits
\newcommand{\ldw}[1]{}                           % NOTHING

%%%%%%%%%%%%%%%%%%%%%%%%%%%%%%%%%%%%%%%%%%%%%%%%%%%%%%%%%%%%%


%%%%%%%%%%%%%%%%%%%%%%%%%%%%%%%%%%%%%%%%%%%%%%%%%%%%%%%%%%%%%%%%%%%%%%%%%%

\title{Aliquot sequences starting with a number under 10000}
\author{
Manuel Benito
\and 
Juan L. Varona
}
\date{}

\begin{document}
\maketitle 

\begin{abstract}
In this paper we describe some advances in the knowledge of the behavior of
aliquot sequences starting with a number less than $10000$. We give the URL address
of the web pages with the free number theory packages that we have used to compute the
aliquot sequences up to their present status.
\end{abstract}

\section*{Introduction and results}

For a positive integer $n$, let $\sigma(n)$ denote the sum of its divisors
(including $1$ and $n$), and $s(n) = \sigma(n) - n$
the sum of its divisors without $n$.
A perfect number is a number $n$ such that
$s(n) = n$, and an amicable pair of numbers is a pair $(n,m)$ satisfying
$s(n) = m$, $s(m) = n$. In a similar way, cycles of numbers
$(a_1, a_2, \dots, a_l)$ such that $s(a_i) = a_{i+1}$ for $1 \le i \le l-1$
and $s(a_l) = a_1$ are known as aliquot cycles or sociable numbers.

The function $s$ generates an aliquot sequence $\{s^k(n)\}_{k=0}^\infty$ 
by taking $s^0(n) = n$ and $s^{k+1}(n) = s(s^k(n))$. For each one of these sequences,
there are four possibilities:
\begin{itemize}
\setlength{\itemsep}{0pt}
\item[$(i)$] it terminates at $1$ (being the previous term a prime number),
\item[$(ii)$] it reaches a perfect number,
\item[$(iii)$] it reaches an amicable pair or a cycle,
\item[$(iv)$] it is unbounded.
\end{itemize}
The Catalan-Dickson conjecture~\cite{Cat,Dick} says that $(iv)$ does not actually happen.
But other researchers disagree with this conjecture and think that there
are unbounded sequences; in fact, the alternative conjecture from
Guy-Selfridge~\cite{Guy-Sel} (see also~\cite{Guy-Guy}, \cite{Guy}) states that
there are many sequences that go to infinity, perhaps almost all of
those that start at an even number (i.e., the proportion of even integers $n$
such that $\{s^k(n)\}_{k=0}^\infty$ is bounded tends to zero).

This alternative conjecture is based upon the existence of several	patterns that, whenever 
they appear in the factor decomposition of $n$, they appear again (with high order of
probability) in the decomposition of $s(n)$. These patterns are called \emph{drivers} or
\emph{guides} (see~\cite{Guy-Sel} for details; also, see an example later in this paper).


Let us explain how to compute a aliquot sequence. It is enough to describe how to compute
$\sigma(n)$ for any positive integer~$n$. Let $n = p_1^{a_1} \cdots p_k^{a_k}$ be its
decomposition in prime numbers. We claim that
\begin{equation}
\label{ec:sigmaSimple}
    \sigma(p_1^{a_1} \cdots p_k^{a_k})
    = (1+p_1+\cdots+p_1^{a_1}) \cdots (1+p_k+\cdots+p_k^{a_k}).
\end{equation}
Indeed, if  we expand the expression on the right, there appear, as summands,
all the divisors of~$p_1^{a_1} \cdots p_k^{a_k}$. Moreover,
$1+p+\cdots+p^a = \frac{p^{a+1}-1}{p-1}$
so we also have the following formula to compute~$\sigma(n)$:
\begin{equation}
\label{ec:sigma}
    \sigma(n) = \frac{p_1^{a_1+1}-1}{p_1-1} \cdots
      \frac{p_k^{a_k+1}-1}{p_k-1}.
\end{equation}

Now, let us suposse $n = 2^3 \cdot 3 \cdot 5 \cdot p \cdot q$, with~$p$ and~$q$
prime numbers, $p \ne q$, $p,q>5$ (the case with more prime factors is similar).
By~(\ref{ec:sigmaSimple}),
$s(n) = (1+2+4+8) \cdot 4 \cdot 6 \cdot (p+1) \cdot (q+1) - 2^3 \cdot 3 \cdot 5
\cdot p \cdot q = 15 \cdot 2^2 \cdot 2 \cdot 3 \cdot (p+1) \cdot (q+1)
- 2^3 \cdot 3 \cdot 5 \cdot p \cdot q = 2^3 \cdot 3 \cdot 5 \cdot
[3 \cdot (p+1) \cdot (q+1) - p \cdot q] = 2^3 \cdot 3 \cdot 5 \cdot m$,
with~$m$ an odd number that is not divisible by~$3$. Then, we get again $2^3 \cdot 3
\cdot 5$ with~$2$ and~$3$ raised to the same power as before.
In this way, we say that $2^3 \cdot 3 \cdot 5$ is a driver. However, the power of $5$ can,
eventually, change; this may cause the breaking of the driver at a later moment.

The smallest $n$ for which there was ever doubt about the end of
$\{s^k(n)\}_{k=0}^\infty$ was $138$, but Lehmer showed that the sequence terminates at
$s^{177}(138) = 1$. Currently, the first number whose behavior is not known is $276$. Many
researchers have investigated the behavior of this sequence, but it is not yet known if its
end exists.

But there are many other sequences whose end is in doubt. For instance, there are five main
sequences starting in a number smaller than $1000$ whose behavior is unknown (Lehmer five:
$276$, $552$, $564$, $660$ and $966$). Before the year $1980$, there were fourteen main
sequences starting between $1000$ and $2000$ (Godwin fourteen), but at the present time only
twelve remain unknown (Godwin showed that the $1848$ sequence terminates, and so did
Dickerman for the $1248$ one).

In 1994, A. W. P. Guy and R. K. Guy published an article~\cite{Guy-Guy} with a table showing
the status of the sequences starting with numbers less than or equal to $7044$. In the same
year, the book~\cite{Guy}, from R.~K.~Guy, updates the information of the article. In this
book, a extensive bibliography on this subject can be found.

Since then, the authors have been checking the aliquot sequences for numbers under $10000$.
For some starting values, we have shown for the first time that the sequence terminates. In
December 22, 1996, we found the record for the maximum of a terminating sequence:
the one starting at $4170$ converges to $1$ after $869$ iterations, getting a maximum
of 84 decimal digits at iteration~$289$. This result was published in~\cite{Ben-Var}.

Later, in October 1999, W.~Bosma broke this record: he found that the aliquot sequence
starting with $44922$ terminates after $1689$ iterations (at 1), reaching a maximum of
$85$ digits at step $1167$. In December 3, 1999, he broke again the record by finding
that the sequence starting at $43230$ finished: it terminates (at $1$) after
$4357$ iterations, reaching a maximum of $91$ digits at step~$967$.

We have continued analyzing the behavior of the aliquot sequences starting with a number
$n < 10000$. In Table~\ref{tabla} we summarize our work with such sequences; it is an update
of the one that appears in~\cite{Ben-Var}. In the table we show the main sequences whose
status is yet unknown (by \textit{main} we mean that all the other sequences with
unknown end merge with any sequence in the table; in this case, we only show
in the table the sequence that starts with the smaller number). We also include the
number of decimal digits of the last term $s^k(n)$ reached for each sequence
and the guide in that stage. Actually, all the sequences are in driver (a driver is more
stable than a guide), except the ones marked with~$^{(*)}$.

In~\cite{Ben-Var}, any of the sequences of the table had been pursued up to, at least, $75$
digits. Now, we have reached more that $92$ digits for all of them; $100$ digits for many of
them. The sequences corresponding $276$, $552$, $564$, $660$, $966$, $1074$, $1134$ are due
to W.~Creyaufm\"uller, P.~Zimmermann and J.~Howell.
Finally, let us note that we have reached at least a hundred digits for the Godwin
sequences.

\begin{table}
\caption{Aliquot sequences whose end is in doubt} \label{tabla}
%\renewcommand{\topfraction}{0.95}
%\renewcommand{\bottomfraction}{0.95}
%\renewcommand{\textfraction}{0.05}
\begin{center}
\renewcommand{\arraystretch}{0.80}
%\small % To choose: \small, \footnotesize or \scriptsize.
\begin{tabular}{|c|ccccccccc|}
%\begin{tabular}{|c|c@{\,}c@{\,}c@{\,}c@{\,}c@{\,}c@{\,}c@{\,}c@{\,}c|}
\hline
\info &
\num{276}{1132}{108}{2 \pp 3} &
\num{552}{818}{118}{2^5 \pp 3 \pp 7} &
\num{564}{2970}{102}{2^2 \pp 7} &
\num{660}{439}{110 \ldw{104}}{2^2 \nodriver} & % was 2^5*3*7 at i371
\num{966}{481}{108}{2 \pp 3} & %%% 2 \pp 3^2
\num{1074}{1543}{103}{2^2 \pp 7} &
\num{1134}{2242}{124}{2^5 \pp 3 \pp 7} & %%% 2^5 \pp 3^4 \pp 7
\num{1464}{1897}{101}{2^2 \pp 7} &
\num{1476}{1055}{106}{2^3 \pp 3 \pp 5}
\\ \hline
\info &
\num{1488}{824}{103}{2^2 \pp 7} &
\num{1512}{1632}{101}{2^2 \pp 7} &
\num{1560}{1336}{101}{2^5 \pp 3 \pp 7} & %%% 2^5 \pp 3^2 \pp 7
\num{1578}{1109}{104 \ldw{88}}{2^2 \nodriver} &
\num{1632}{713}{102}{2^3 \pp 3 \pp 5} & %%% 2^3 \pp 3^2 \pp 5
\num{1734}{1393}{101}{2^2 \pp 7} &
\num{1920}{1992}{108}{2^2 \pp 7} & %%%% CAUSALIDAD: LA SIGUIENTE SUCESION ES LA 1992
\num{1992}{985}{102}{2^3 \pp 3 \pp 5} & %%% 2^3 \pp 3^3 \pp 5
\num{2232}{390}{102 \ldw{95}}{2^6 \nodriver}
\\ \hline
\info &
\num{2340}{471}{99}{2^3 \pp 3 \pp 5} &
\num{2360}{974}{95}{2^2 \pp 7} &
\num{2484}{796}{97}{2^3 \pp 3} & %%% 2^3 \pp 3^3
\num{2514}{2836}{95}{2^3 \pp 3 \pp 5} & %%% 2^3 \pp 3 \pp 5^2
\num{2664}{761}{100}{2^2 \pp 7} &
\num{2712}{1347}{95}{2^2 \pp 7} &
\num{2982}{810}{97}{2^4 \pp 31} &
\num{3270}{417}{98}{2^5 \pp 3 \pp 7} & %%% 2^5 \pp 3 \pp 7^2
\num{3366}{1062}{100 \ldw{82}}{2^3 \nodriver}
\\ \hline
\info &
\num{3408}{840}{95}{2^3 \pp 3 \pp 5} &
\num{3432}{933}{103}{2^3 \pp 3 \pp 5} &
\num{3564}{779}{100}{2^3 \pp 3} & %%% 2^3 \pp 3^2
\num{3630}{1240}{100 \ldw{94}}{2^2 \nodriver} &
\num{3678}{1201}{98}{2^2 \pp 7} &
\num{3774}{1193}{98 \ldw{89}}{2^7 \pp 3 \nodriver} &
\num{3876}{830}{96}{2^2 \pp 7} &
\num{3906}{679}{94}{2^2 \pp 7} &
\num{4116}{1192}{105}{2^3 \pp 3}
\\ \hline
\info &
\num{4224}{519}{98}{2^3 \pp 3 \pp 5} &
\num{4290}{913}{92}{2^2 \pp 7} &
\num{4350}{1165}{97}{2^2 \pp 7} &
\num{4380}{965}{100}{2^2 \pp 7} & %%% 2^2 \pp 7^2
\num{4788}{2152}{105}{2^3 \pp 3 \pp 5} & %%% 2^3 \pp 3^3 \pp 5
\num{4800}{1135}{101}{2^4 \pp 31} &
\num{4842}{473}{98}{2^2 \pp 7} &
\num{5148}{1545}{95}{2 \pp 3} &
\num{5208}{1710}{96}{2^3 \pp 3 \pp 5}
\\ \hline
\info &
\num{5250}{1564}{100 \ldw{89}}{2^4 \nodriver} &
\num{5352}{683}{93}{2^2 \pp 7} &
\num{5400}{2696}{93}{2^2 \pp 7} &
\num{5448}{1185}{96}{2^3 \pp 3 \pp 5} &
\num{5736}{1093}{100}{2^2 \pp 7} &
\num{5748}{1045}{98}{2^2 \pp 7} &
\num{5778}{742}{95}{2^4 \pp 31} &
\num{6160}{1630}{96 \ldw{84}}{2^2 \nodriver} &
\num{6396}{1234}{92}{2^3 \pp 3 \pp 5} %%% 2^3 \pp 3 \pp 5^2
\\ \hline
\info &
\num{6552}{893}{93}{2^3 \pp 3} &
\num{6680}{1815}{94}{2^4 \pp 31} &
\num{6822}{1177}{97}{2^4 \pp 31} &
\num{6832}{885}{104}{2^3 \pp 3} &
\num{6984}{1764}{96}{2^4 \pp 31} &
\num{7044}{1113}{102}{2^4 \pp 31} &
\num{7392}{498}{96}{2^3 \pp 3 \pp 5} & %%% 2^3 \pp 3^2 \pp 5^2
\num{7560}{846}{97 \ldw{86}}{2^3 \pp 5 \nodriver} &
\num{7890}{891}{99}{2^2 \pp 7} %%% 2^2 \pp 7^2
\\ \hline
\info &
\num{7920}{951}{95}{2^6 \pp 127} &
\num{8040}{2205}{94}{2^3 \pp 3 \pp 5} &
\num{8154}{647}{96}{2^2 \pp 7} &
\num{8184}{1241}{102}{2^2 \pp 7} &
\num{8288}{849}{103 \ldw{95}}{2^8 \nodriver} &
\num{8352}{1291}{96}{2^2 \pp 7} &
\num{8760}{2145}{94}{2^3 \pp 3 \pp 5} & %%% 2^3 \pp 3^3 \pp 5^3
\num{8844}{1184}{101}{2^4 \pp 31} &
\num{8904}{963}{95}{2^2 \pp 7}
\\ \hline
\info &
\num{9120}{532}{92}{2^3 \pp 3} &
\num{9282}{505}{95}{2^3 \pp 3} &
\num{9336}{608}{97}{2^6 \pp 127} &
\num{9378}{2111}{92}{2^2 \pp 7} &
\num{9436}{545}{94}{2 \pp 3} &
\num{9462}{447}{97}{2^3 \pp 3 \pp 5} &
\num{9480}{945}{93}{2 \pp 3} &
\num{9588}{1848}{103}{2^5 \pp 3 \pp 7} &
\num{9684}{617}{92}{2^5 \pp 3 \pp 7}
\\ \hline
\info &
\num{9708}{671}{95}{2^2 \pp 7} &
\num{9852}{669}{105}{2^2 \pp 7} &&&&&&&
\\ \hline
\end{tabular}
\end{center}
\end{table}
%%%%%%%%%%%%%%%%%%%%%%%%%%%%%%%%%%%%%%%%%%%%%%%%%

\section*{The method}

According to~(\ref{ec:sigmaSimple}) or~(\ref{ec:sigma}), to compute an aliquot sequence, we
have to decompose $m$ in factors to calculate~$s(m)$.
This is hard when the number $m$ is big. And we are dealing with numbers more
than a hundred of digits long, so powerful algorithms are necessary. In particular, we
have used the elliptic curve method (ECM) and the multipolynomial quadratic sieve (MPQS). We
have checked the primality of the factors with the Adleman-Pomerance-Rumely (APR) primality
test.

All our work has been done by using free packages available on internet.
We have run the programs on many computers from the authors and some colleages, and their
respective institutions.
% The work has been done mostly during nights and weekends.

We have used the following packages, that are available at their corresponding web pages
(or anonymous \texttt{ftp} sites):

\begin{itemize}
\setlength{\itemsep}{0pt}

\item UBASIC,
\url{ftp://rkmath.rikkyo.ac.jp/pub/ubibm/}

\item PARI-GP,
\url{http://www.parigp-home.de/}

\item KANT-KASH,
\url{http://www.math.TU-Berlin.DE/algebra/}

\item MIRACL,
\url{http://indigo.ie/~mscott/}

\end{itemize}


\section*{Other people and the web}

There are some other people working in aliquot sequences:
Wolfgang Creyaufm\"uller is computing aliquot sequences starting with a number between
$10^5$ and $10^6$ up to a minimum of $60$ digits; also, he has written the
book~\cite{Crey}.
Paul Zimmermann is computing iterates of the sequences starting in the numbers $276$,
$552$, $564$, $660$, $966$, $1074$, $1134$ and $204828$; also, he is
working with aliquot sequences starting between $50000$ and~$100000$.
Jim Howell is working in the aliquot sequence starting in~$966$.
Finally, Wieb Bosma is computing with aliquot sequences starting between $10000$ and~$50000$.

One of the authors maintains the following web related with aliquot sequences:
% LO DEL \quad ES UN TRUCO PARA QUE SE VEA BIEN EL ENLACE EN EL .pdf.
\begin{center}
\quad\url{http://www.unirioja.es/dptos/dmc/jvarona/aliquot.html}\quad
\end{center}
In this page, you can find links to the aliquot web pages of
W.~Creyaufm\"uller, P.~Zimmermann, J.~Howell and W.~Bosma. Also, the
sequences themselves can be found in the \texttt{ftp} server
\begin{center}
\quad\url{ftp://mat.unirioja.es/pub/aliquot/00xxxx}\quad
\end{center}

\section*{Acknowledgments}

The authors thank Wolfgang Creyaufm\"uller, Paul Zimmermann and Jim Howell for their data on
the sequences starting in the numbers $276$, $552$, $564$, $660$, $966$, $1074$, $1134$,
which can be found in their web pages.

Finally, thanks to Laureano Lamb\'an, Eloy Mata, David Ortigosa, Julio Rubio
and Divina S\'aenz for allowing us to use their personal computers to do some of the
computations that appear in this note.



%%%%%%%%%%%%%%%%%%%%%%%%%%%%%%%%%%%%%%%%%%%%%%%%%%%%%%%%%%%%%%%%%%%%%%%%%%%
\begin{thebibliography}{9}
%%%%%%%%%%%%%%%%%%%%%%%%%%%%%%%%%%%%%%%%%%%%%%%%%%%%%%%%%%%%%%%%%%%%%%%%%%%

\bibitem{Ben-Var} M. Benito and J. L. Varona,
Advances in aliquot sequences,
\textit{Math. Comp.}
\textbf{68} (1999), 389--393.

\bibitem{Cat} E. Catalan,
Propositions et questions diverses,
\textit{Bull. Soc. Math. France}
\textbf{18} (1887--88), 128--129.

% \bibitem{Cohen} H. Cohen,
% \textit{A course in computational algebraic number theory},
% Springer-Verlag, 1993.

\bibitem{Crey} W. Creyaufm\"uller,
\textit{Primzahlfamilien} (3\textsuperscript{rd}~ed.),
Verlagsbuchhandlung Creyaufm{\"u}ller, Stuttgart, 2000.
% En alem\'an, 317 p., ISBN 3-9801032-2-6

\bibitem{Dick} L. E. Dickson,
Theorems and tables on the sum of the divisors of a number,
\textit{Quart. J. Math.}
\textbf{44} (1913), 264--296.

\bibitem{Guy-Guy} A. W. P. Guy and R. K. Guy,
A record aliquot sequence,
in \textit{Computation 1943--1993: A Half-Century of
Computational Mathematics} (Vancouver, 1993),
\textit{Proc. Sympos. Appl. Math.}
\textbf{48} (1994),
Amer. Math. Soc., Providence, RI, 1994,
pp.~557--559.
% %\MR{95k:11164}

\bibitem{Guy} R. K. Guy,
\textit{Unsolved Problems in Number Theory} (2\textsuperscript{nd}~ed.),
Springer-Verlag, 1994.
%\MR{96e:11002}

\bibitem{Guy-Sel} R. K. Guy and J. L. Selfridge,
What drives an aliquot sequence?,
\textit{Math. Comp.}
\textbf{29} (1975), 101--107.
Corrigendum, \textit{ibid.} \textbf{34} (1980), 319--321.
%\MR{52:5542}; \MR{81f:10008}

\end{thebibliography}

%%%%%%%%%%%%%%%%%%%%%%%%%%%%%%%%%%%%%%%%%%%%%%%%%%%%%%%%%%%%%%%%%%%%%%%%%%%
%%%% DIRECCION DE LOS AUTORES DEL ART\'ICULO
%%%%%%%%%%%%%%%%%%%%%%%%%%%%%%%%%%%%%%%%%%%%%%%%%%%%%%%%%%%%%%%%%%%%%%%%%%%

\vskip 6pt

\begin{flushleft} % O flushright, SWGUN SE QUIERA
\noindent
{\small Manuel Benito,\\
Instituto Sagasta,\\
C/ Doctor Zub\'{\i}a s/n, 26003 Logro\~no.\\
e-mail: \href{mailto:mbenit8@palmera.pntic.mec.es}{mbenit8@palmera.pntic.mec.es}
}
\end{flushleft}

\vskip 6pt

\begin{flushleft}
{\small Juan Luis Varona,\\
Departamento de Matem\'aticas y Computaci\'on,
Universidad de La Rioja,\\
C/ Luis de Ulloa s/n, 26004 Logro\~no.\\
e-mail: \href{mailto:jvarona@dmc.unirioja.es}{jvarona@dmc.unirioja.es}\\
URL: \url{http://www.unirioja.es/dptos/dmc/jvarona/welcome.html}
\phantom{URL: }\url{http://www.unirioja.es/dptos/dmc/jvarona/hola.html}
}
\end{flushleft}

\end{document}

